\subsection{Computational complexity by workload}
\label{sec:sncs-complexity}
Our comparison separates the cost of (i) preprocessing,
(ii) one complete trend, (iii) one criterion score,
(iv) a search over $M$ candidate penalties, and
(v) repeatedly evaluating the same candidates for $P$
criteria or $R$ independent series. These workloads
are \emph{not interchangeable}. Let $n$ denote the
number of observations and $d$ the fixed difference
order. The normal-equation matrix
$A_\lambda=I+\lambda D_d^\top D_d$
has half-bandwidth $d$. Efficient structured
Whittaker--Henderson solutions substantially predate
our implementation \cite{weinert2007wh}.

\begin{table}[tbp]
\centering\scriptsize\setlength{\tabcolsep}{3pt}
\caption{Leading operation and memory orders for
different numerical solvers. A new penalty entails
a new factorization for the direct methods, whereas
spectral bases can be reused. CV leverage and
trace calculations may add substantial work.
The expressions assume $d\ll n$; these are
complexities, not measured runtime rankings.}
\label{tab:sncs-solver-complexity}
\begin{tabular}{p{.29\linewidth}p{.23\linewidth}p{.24\linewidth}p{.12\linewidth}}
\toprule
Method & Setup & Full trend per $\lambda$ & Memory\\
\midrule
Dense solve or dense Cholesky
& $O(nd^2)$ for $Q_d$
& $O(n^3)$ factorization
& $O(n^2)$\\
General banded solve
& Store $2d+1$ diagonals
& $O(nd^2)$
& $O(nd)$\\
Banded Cholesky
& Store $d+1$ diagonals
& $O(nd^2)$ factor; $O(nd)$ solve
& $O(nd)$\\
SVD of $D_d$
& $O(n^3)$
& $O(n^2)$
& $O(n^2)$\\
Symmetric \texttt{eigh} of $Q_d$
& $O(n^3)$
& $O(n^2)$
& $O(n^2)$\\
\bottomrule
\end{tabular}
\end{table}

After dense spectral preprocessing, an EDF/trace
evaluation and a scalar normalized $S(\lambda)$
evaluation require $O(n)$ operations. Once the
observations have been projected, spectral
GCV, AICc, BIC and the implemented conditional
Gaussian likelihood also admit $O(n)$ score
paths per penalty. Full trend reconstruction
and the usual spectral diagonal-leverage
computation for ordinary CV require
$O(n^2)$ per candidate. These costs explain
why a fast one-curve banded solver need not
be the fastest method for a complete search.

For $M$ distinct \emph{full fitted curves},
two leading cost models are
\[
 T_{\rm spectral}=O(n^3+Mn^2),\qquad
 T_{\rm banded\ Cholesky}=O(Mnd^2).
\]
The spectral cost includes its one-time basis.
Banded Cholesky exploits fixed bandwidth
but must refactor for each different $\lambda$.
The constants, BLAS implementation, order
and requested diagnostic quantities determine
the actual crossover.

\subsection{The cost of an $S$-only coordinate}
\label{sec:sncs-mapping-cost}
The normalized mapping $S_d(\lambda)$ depends
only on $n,d$ and the positive penalty spectrum;
it does not depend on $y$ or on the fixed
reference difference mean $\mu$.
After preprocessing, one evaluation of $S$
or its derivative costs $O(n)$.
A safeguarded inversion with $k$ scalar
iterations costs $O(kn)$ per target.
A full $M$-target grid therefore requires
$O(Mkn)$ arithmetic \emph{in addition to}
the spectral setup. If a banded Cholesky
backend does not already have the spectrum,
obtaining the mapping has a separate setup
cost; the banded factorization is not itself
a free $S^{-1}$ routine.

Suppose $P$ criteria use the same $M$ finite
$S$ targets. Writing $T_{\rm inv}$ for
a candidate inversion and $T_{\rm score}$
for a typical matched criterion evaluation,
the two strategies satisfy the idealized
accounting identity
\begin{equation}
\begin{aligned}
T_{\rm native}&\approx T_{\rm prep}
  +PM(T_{\rm inv}+T_{\rm score}),\\
T_{\rm cached}&\approx T_{\rm prep}
  +MT_{\rm inv}+PMT_{\rm score},\\
T_{\rm native}-T_{\rm cached}
  &\approx (P-1)MT_{\rm inv}.
\end{aligned}
\label{eq:sncs-mapping-reuse}
\end{equation}
For unequal criterion costs, replace the
common score term by the sum over criteria.
This relationship assumes a \emph{matched}
candidate set and a one-time cached
translation shared among all criteria.
At $P=1$ the conversion cannot generally
be eliminated: it is merely moved to an
earlier stage. For repeated series with the
same $(n,d)$ and target grid, the cached
mapping remains reusable.

By contrast, searching in $\lambda$
and reporting $S(\widehat\lambda)$
at the end avoids repeated inverse
calculations. It can be cheaper for
one-off selection. However, sampling a
uniform $S$ grid and a log-$\lambda$
grid produces \emph{different} interior
candidates; any different discrete
minimum may arise from grid design.
Exact monotonic reparameterization
preserves the global mathematical
minimizer on the full domain.

\subsection{Measured solver crossovers and grid reuse}
\label{sec:sncs-measured-crossovers}
Our paired full-trend results
(Table~\ref{tab:expanded-solver-orders})
show that no backend universally wins.
For $n=32,d=2,\mu=0$, the SVD/banded
Cholesky total-time ratio was $1.554$
for $M=1$ (pointwise 95\% interval
$[1.541,1.568]$), favoring Cholesky,
but $0.419$ for $M=64$
($[0.417,0.421]$), favoring SVD.
At $n=96,d=2,M=16,\mu=0$,
the ratio was $2.093$
($[2.084,2.101]$), again favoring
banded Cholesky. These host-specific
measurements do not justify a
universal crossover threshold.

An independent Windows result
(Table~\ref{tab:expanded-coordinates})
shows the value of reusing the
\emph{same} translated grid across
multiple statistical selectors:
$n=48,d=2,\mu=0$ required
$11.234$\,ms for five repeated native
$S$ searches, compared with
$3.980$\,ms for one $S$-to-$\lambda$
translation plus five backend
searches. The ratio $2.823$ is
a descriptive measured advantage
of avoiding redundant inversions,
not proof that $S$ is intrinsically
faster than $\lambda$. An
independent log-$\lambda$ grid
is not an exact matched comparison.

\subsection{A public interface that exposes $S$, not $\lambda$}
\label{sec:sncs-s-only-api}
The package can make normalized $S$
the only \emph{required user-specified}
smoothing parameter while retaining
$\lambda$ internally as a numerical
representation. The new additive
\texttt{fit\_smoothness} interface
calls the existing, validated
$S^{-1}$ algorithm and then the
usual PLS solver:
\begin{verbatim}
from trendestimation import fit_smoothness

result = fit_smoothness(
    y, smoothness=0.86, order=2, mu=0.15,
)
trend = result.trend
\end{verbatim}
The returned penalty remains
available as optional reproducibility
metadata. Reusing
\texttt{TrendSmoother.fit\_smoothness}
avoids repeating spectral setup
between calls, although repeated
grid inversions should themselves
be cached when appropriate.
The familiar $\lambda$-based API
remains supported for compatibility.
Thus eliminating $\lambda$ from
the interface does \emph{not}
eliminate its internal computation
or constitute a new estimator.
